\documentclass[journal]{IEEEtran}
\usepackage{microtype}
\usepackage{amsmath}
\usepackage{booktabs}

\begin{document}

\title{Paragraph-Local Recompilation for Interactive
Editing of \LaTeX{} Journal Manuscripts}

\author{A.~Author,~B.~Author,~and~C.~Author%
\thanks{The authors are with the Institute for Typesetting Research,
Chennai, India (e-mail: authors@typesetting-research.example).}%
\thanks{Manuscript received July 21, 2026.}}

\markboth{Journal of Document Engineering, Vol.~1, No.~1, July~2026}%
{Author \MakeLowercase{\textit{et al.}}: Paragraph-Local Recompilation}

\maketitle

\begin{abstract}
Interactive editing of large \LaTeX{} manuscripts is limited by
whole-document compilation, whose latency grows with document length and
preamble complexity. We show that the unit of recompilation can instead be
the paragraph: line breaking is a pure function of paragraph content and a
finite parameter context, so an edited paragraph can be re-typeset in
isolation in one to three milliseconds inside a persistent engine, while a
background pass reconverges global layout. We describe the required context
capture, measure fidelity and latency on production journal classes, and
demonstrate a browser-based editor whose keystroke path runs two orders of
magnitude inside the interactive frame budget.
\end{abstract}

\begin{IEEEkeywords}
Typesetting, line breaking, real-time systems, document engineering.
\end{IEEEkeywords}

\section{Introduction}

\IEEEPARstart{A}{utomated} typesetting systems have historically treated
document compilation as a monolithic operation: source files enter, a
complete rendered document leaves. This architecture, inherited from the
batch environments of the nineteen-seventies, imposes a latency floor on
interactive editing that grows linearly with document length and with the
severity of the preamble.

The line-breaking problem, by contrast, is intrinsically local. Given a
fixed measure and a fixed sequence of boxes, glue, and penalties, the
optimal arrangement of breakpoints depends on nothing outside the paragraph
itself. This observation, implicit in the dynamic programming formulation
of the classical algorithm, suggests that the unit of interactive
recompilation should be the paragraph rather than the document.

Word processors have exploited exactly this separation for decades: the
paragraph under the cursor is re-flowed on every keystroke, while
pagination, running heads, and cross-references converge in a background
process. What has been missing is not an algorithm but an engineering
bridge that brings the same interaction model to a batch typesetting
engine without sacrificing its output quality.

This paper contributes such a bridge. We capture, for every paragraph of a
real journal manuscript, the complete parameter context consumed by the
line breaker; we keep a persistent engine alive so that the preamble is
paid for once; and we verify, glyph by glyph, that isolated recompilation
reproduces the reference document exactly.

\section{The Locality Property}

Consider a paragraph consisting of $n$ items with natural widths $w_i$,
stretchability $y_i$, and shrinkability $z_i$. The adjustment ratio of a
candidate line is $r = (L - W)/Y$ when the line must stretch, where $L$ is
the target measure, $W$ the accumulated natural width, and $Y$ the
accumulated stretchability. The badness of setting the line is
approximately $100\,|r|^3$, capped at ten thousand.

The total demerits of a breakpoint sequence combine badness with penalty
terms,
\begin{equation}
D = \sum_{j} \bigl( l + \beta_j \bigr)^2 + \pi_j^2 + \delta_j,
\label{eq:demerits}
\end{equation}
where $l$ is the line penalty, $\beta_j$ the badness of line $j$, $\pi_j$
any explicit penalty at the chosen breakpoint, and $\delta_j$ the adjacency
demerits charged when consecutive lines differ excessively in tightness.

Equation~\eqref{eq:demerits} is minimized by dynamic programming over
feasible breakpoints, and the minimizing sequence is unique for practical
parameter settings. The consequence relevant to interactive systems is
determinism: recompiling an unchanged paragraph reproduces byte-identical
line boxes, and recompiling an edited paragraph perturbs no other
paragraph in the document.

The dependency set of the line breaker is finite and enumerable: the
measure, the paragraph shape, the fonts active at paragraph start, the
hyphenation language, the tolerance and emergency stretch, the penalty and
demerit parameters, and the interword skip overrides. We refer to this
collection as the paragraph's \emph{context}. If the context is captured
during a reference compilation and reinjected around an isolated
recompilation, the reproduction is exact.

\section{System Architecture}

Our prototype comprises four cooperating components.

First, a capture pass instruments a normal compilation. Callbacks around
the line breaker record each paragraph's context together with a reference
signature: the width, height, and depth of every line box and the exact
position of every glyph within it, expressed in scaled points, of which
there are $65{,}536$ to the printer's point\footnote{All measurements in
this paper were taken with LuaHBTeX 1.22.0 from \TeX{} Live 2025 on a
commodity Linux workstation.}.

Second, a persistent engine loads the journal class preamble once and then
services recompilation requests indefinitely. Each request supplies the
edited source of one paragraph together with its captured context, which
the server reinstates before re-running the paragraph builder; the
response carries positioned glyphs extracted directly from the engine's
node structures, bypassing page generation entirely.

Third, cross-reference state is transplanted from the reference
compilation's auxiliary file, so that citations and references inside an
isolated paragraph resolve to their converged values rather than to
placeholder marks.

Fourth, a background convergence pass recompiles the full document a short
interval after the last edit, refreshing the cached page geometry against
which live paragraphs are overlaid. Pagination, floats, and running heads
are therefore eventually consistent, in the same sense a word processor's
layout is.

\section{Evaluation}

We evaluate fidelity and latency on the journal class of this very
manuscript. Fidelity is judged by the strictest available criterion: the
recompiled paragraph must reproduce the reference line breaks and every
glyph position to zero scaled points, with full character protrusion and
font expansion active.

Across the body paragraphs of the test manuscript, isolated recompilation
is exact wherever the paragraph does not couple to global state. The
exceptions are principled rather than accidental: a footnote couples its
paragraph to the page builder through the insertion mechanism, and a
displayed equation splits its paragraph into partial paragraphs whose
shapes depend on the display. Both cases are detectable by inspection of
the node list and are routed to the convergence path.

Latency is summarized in Table~\ref{tab:latency}. The keystroke path
completes in single-digit milliseconds, roughly two orders of magnitude
inside the hundred-millisecond threshold at which interfaces are judged
instantaneous, and comfortably inside a sixty-hertz display frame.

\begin{table}[!t]
\caption{Measured Latency of the Keystroke Path\label{tab:latency}}
\centering
\begin{tabular}{lcc}
\toprule
Stage & Short para. & Medium para. \\
\midrule
Line break + traversal & 0.3\,ms & 2.0\,ms \\
Interprocess transport & 0.2\,ms & 0.5\,ms \\
Full round trip & 0.5\,ms & 2.5\,ms \\
Full document compile & \multicolumn{2}{c}{900\,ms} \\
\bottomrule
\end{tabular}
\end{table}

The comparison with whole-document recompilation is stark. Even this
modest manuscript costs nearly a second per compile; a book-length
document costs tens of seconds. The per-paragraph path is independent of
document length by construction, so the advantage widens without bound as
the document grows.

Memory behavior is equally important for a long-lived engine. Every
recompiled paragraph allocates node lists that must be freed
deterministically; our server frees the working box after each response,
and five hundred consecutive recompilations show no drift in either
latency or engine memory.

\section{Limitations and Future Work}

Three limitations bound the present prototype. Paragraphs inside list
environments receive their labels from the enclosing environment rather
than from their own source, so label glyphs are absent from isolated
recompiles. Global assignments in body text can leak between requests in a
long-lived engine; checkpointing would make isolation airtight. And
mid-edit source is frequently ill-formed, which the server survives by
sanitizing input and respawning on the rare wedge, but a production system
would maintain an incremental parse of the manuscript instead.

The natural next step is source mapping: deriving paragraph identity from
an editor's document model rather than from markers, so that any existing
manuscript can be edited live without preparation. The capture and replay
machinery described here is unchanged by that step.

\section{Conclusion}

Pixel-identical \LaTeX{} output and millisecond editing feedback are not
in tension. The line breaker that produces the former has always been fast
enough for the latter; what was required was to stop asking it to
recompile entire documents. A paragraph-local fast path with background
convergence brings the interaction model of word processors to batch
typesetting without compromising its typography.

\begin{thebibliography}{9}
\bibitem{knuthplass} D.~E. Knuth and M.~F. Plass, ``Breaking paragraphs
into lines,'' \emph{Software---Practice and Experience}, vol.~11, no.~11,
pp. 1119--1184, 1981.
\bibitem{lode} C.~Lode, ``Real-time Lua\TeX: recompiling large documents
in 1\,ms,'' \emph{TUGboat}, 2026, preprint.
\bibitem{bour} F.~Bour, ``TeXpresso: live rendering and error reporting
for \LaTeX,'' \emph{TUGboat}, vol.~44, no.~2, pp. 185--192, 2023.
\end{thebibliography}

\end{document}
